%% Use the RPI thesis document class.
\documentclass[chap]{rpi_thesis}

% Provides 1 in. margins on all sides of the page.
% Adding the option 'showframe' to the \geometry options below
% will show an outline of the margins on all pages.
\usepackage{geometry}
\geometry{
left=1in,
right=1in,
top=1in,
bottom=1in}

% These options will provide most of the options for an IEEE
% compliant bibliography and captions for figures. It is up to users to
% capitilizations / abbreviations and other formatting is correct
% in the `references.bib' file.
\usepackage{cite}
\usepackage{notoccite}
\bibliographystyle{rpi_IEEE}
\renewcommand{\bibname}{}

% Consistently caption figures and tables  with a period
% separating the figure number and the figure caption and
% comply with the IEEE style guide for figure options.
\usepackage[
  labelsep=period,
  justification=centering]{caption}
\renewcommand{\captionfont}{\bfseries}
\renewcommand{\figurename}{Fig.}

% Place captions above tables as required by OGE
\usepackage{float}
\floatstyle{plaintop}
\restylefloat{table}

% This package allows for the easy inclusion of figures
\usepackage{graphicx}

% This package allows for hyperrefs in the document, which can
% allow easy navigation to citations and referenced equations,
% chapters, and sections. This is very useful when writing.
\usepackage[hidelinks]{hyperref}

% The following packages are used to generate generic text and
% can be removed once editing begins.
\usepackage[english]{babel}
\usepackage{blindtext}

% The main body of the thesis document. Add or remove files as needed.
\begin{document}
\thesistitle{\bf MY GREAT THESIS}

\author{A Great Graduate Student}
\degree{Doctor of Philosophy}
\department{Mechanical, Aerospace, and Nuclear Engineering}
\submitdate{[May 2018] \\ Submitted January 2018}
\copyrightyear{2018}

\titlepage
\copyrightpage
\tableofcontents
\listoftables
\listoffigures

\specialhead{ACKNOWLEDGMENT}

Thanks to everyone. I'll buy you all ice cream.

\specialhead{ABSTRACT}

\blindtext

%%% INTRODUCTION AND BACKGROUND
\chapter{MY SUPER GREAT CHAPTER}
\label{chap:intro}

\let\thefootnote\relax\footnotetext{
Portions of this chapter have been submitted to:
A.~G. Author, ``A super cool paper,''
submitted for publication.}

%%% INTRODUCTION
\section{Introduction}

All work and no play make Jack a dull boy.
All work and no play make Jack a dull boy.
All work and no play make Jack a dull boy.
All work and no play make Jack a dull boy.
All work and no play make Jack a dull boy.
All work and no play make Jack a dull boy.
All work and no play make Jack a dull boy.
All work and no play make Jack a dull boy.
All work and no play make Jack a dull boy.
All work and no play make Jack a dull boy.
All work and no play make Jack a dull boy.
All work and no play make Jack a dull boy.
All work and no play make Jack a dull boy.
All work and no play make Jack a dull boy.
All work and no play make Jack a dull boy.
All work and no play make Jack a dull boy.
All work and no play make Jack a dull boy.
All work and no play make Jack a dull boy.

\begin{table}[hbt!]
\centering
\begin{tabular}{|c | c |}
\hline
a & b \\ \hline\hline
0.01 & 0.02 \\ \hline
0.03 & 0.04 \\ \hline
\end{tabular}
\caption{An example of a table in the thesis document.
Notice the caption is above the table.}
\label{table:mytable}
\end{table}

All work and no play make Jack a dull boy.
All work and no play make Jack a dull boy.
All work and no play make Jack a dull boy.
All work and no play make Jack a dull boy.
All work and no play make Jack a dull boy.
All work and no play make Jack a dull boy.
All work and no play make Jack a dull boy.

\begin{figure}[ht!]
\centering
\includegraphics[width=0.25\textwidth]{rpi_seal}
\caption{An example of a figure in the thesis document.
Notice the caption is below the figure.}
\label{fig:myfigure}
\end{figure}

Look! I am referencing Figure \ref{fig:myfigure}.
Look! I am referencing Table \ref{table:mytable}.
Look! I am referencing a citation \cite{areference}.

\section{Conclusions}

All work and no play make Jack a dull boy.
All work and no play make Jack a dull boy.
All work and no play make Jack a dull boy.
All work and no play make Jack a dull boy.
All work and no play make Jack a dull boy.
All work and no play make Jack a dull boy.
All work and no play make Jack a dull boy.
All work and no play make Jack a dull boy.
All work and no play make Jack a dull boy.
All work and no play make Jack a dull boy.
All work and no play make Jack a dull boy.
All work and no play make Jack a dull boy.
All work and no play make Jack a dull boy.
All work and no play make Jack a dull boy.
All work and no play make Jack a dull boy.
All work and no play make Jack a dull boy.
All work and no play make Jack a dull boy.
All work and no play make Jack a dull boy.

All work and no play make Jack a dull boy.
All work and no play make Jack a dull boy.
All work and no play make Jack a dull boy.
All work and no play make Jack a dull boy.
All work and no play make Jack a dull boy.
All work and no play make Jack a dull boy.
All work and no play make Jack a dull boy.
All work and no play make Jack a dull boy.
All work and no play make Jack a dull boy.

\specialhead{REFERENCES}
\begin{singlespace}
\bibliography{references}
\end{singlespace}

\appendix
\addtocontents{toc}{\parindent0pt\vskip12pt APPENDICES}

\chapter{MY GREAT ADDITIONAL STUFF}

\section{More Stuff}

All work and no play still make Jack a dull boy.
All work and no play still make Jack a dull boy.
All work and no play still make Jack a dull boy.
All work and no play still make Jack a dull boy.
All work and no play still make Jack a dull boy.
All work and no play still make Jack a dull boy.
All work and no play still make Jack a dull boy.
All work and no play still make Jack a dull boy.
All work and no play still make Jack a dull boy.
All work and no play still make Jack a dull boy.
All work and no play still make Jack a dull boy.

\end{document}
